\title{Topology-Aware Multi-Agent Reinforcement Learning for Joint Resource Allocation in Next-Generation Vehicular Networks}

\begin{abstract}
Next-generation urban vehicular networks must simultaneously satisfy stringent and tightly coupled resource demands across spectrum allocation, energy consumption, and computation offloading. Existing deterministic optimization techniques and flat multi-agent reinforcement learning (MARL) approaches struggle to scale with high vehicle density and fail to generalize across highly dynamic spatial topologies. To address this gap, we propose a novel Hierarchical Graph Attention Multi-Agent Soft Actor-Critic (HiGAT-MASAC) framework to tackle the joint spectrum, power, and Multi-Access Edge Computing (MEC) capacity allocation problem in Beyond 5G (B5G) smart mobility environments. Our architecture decomposes the intractable mixed-integer nonlinear programming problem into a two-timescale Decentralized Partially Observable Markov Decision Process (Dec-POMDP) without sacrificing coordination. At the macro-level, Roadside Unit (RSU) agents periodically allocate globally constrained spectrum sub-bands and computational server capacities; concurrently at the micro-level, distributed vehicle agents rapidly optimize their transmit power and offloading ratios. To inherently capture dynamic Vehicle-to-Vehicle (V2V) interference interactions, we employ Multi-Head Graph Attention Networks (GAT) that natively encode spatial connectivity mapping into expressive state embeddings for the maximum-entropy Soft Actor-Critic (SAC) policies. Through extensive simulations leveraging a structured urban tracking scenario modeled precisely mathematically over the 3GPP channel guidelines, we validate our algorithm’s stability and convergence superiority. Compared to state-of-the-art multi-agent baselines encompassing MAPPO, MADDPG, and GNN-DDQN, the proposed topology-aware hierarchical framework substantially reduces overall task completion delays while maximizing Vehicle-to-Infrastructure (V2I) throughput and operating under sustainable energy bounds. By uniquely disentangling the temporal bandwidth dimension with the continuous edge-offload topology relationships via graph-learned heuristics, this research provides a scalable AI-driven autonomous management system tailored definitively for tomorrow’s urban traffic grids.
\end{abstract}

\begin{keyword}
Multi-Agent Reinforcement Learning \sep Graph Neural Networks \sep Vehicular Networks \sep Resource Allocation \sep Edge Computing \sep Smart Mobility
\end{keyword}
